\clearpage
\subsection{MSVC: x86 + \olly}

\RU{Попробуем в \olly немного хакнуть программу и сделать вид, что \scanf срабатывает всегда без ошибок.}
\EN{Let's try to hack our program in \olly, forcing it to think \scanf always works without error.}
\ES{Intentemos hackear un poco el programa en \olly, obligándolo a pensar que \scanf siempre funciona sin errores.}

\RU{Когда в \scanf передается адрес локальной переменной, изначально в этой переменной
находится некий мусор. В данном случае это}
\EN{When an address of a local variable is passed into \scanf,
the variable initially contains some random garbage, in this case}
\ES{Cuando se pasa la dirección de una variable local a \scanf, la variable contiene inicialmente algo de basura aleatoria, en este caso} \TT{0x6E494714}:

\begin{figure}[H]
\centering
\includegraphics[scale=\FigScale]{patterns/04_scanf/3_checking_retval/olly_1.png}
\caption{\olly: \RU{передача адреса переменной в}\EN{passing variable address into}\ES{pasando la dirección de la variable a} \scanf}
\label{fig:scanf_ex3_olly_1}
\end{figure}

\clearpage
\RU{Когда}\EN{While}\ES{Mientras se ejecuta} \scanf 
\RU{запускается, вводим в консоли что-то непохожее на число, например}
\EN{executes, in the console we enter something that is definitely not a number, like} \q{asdasd}.
\RU{scanf заканчивается с 0 в}\EN{scanf finishes with 0 in} \EAX, \RU{что означает, что произошла ошибка}%
\EN{which indicates that an error has occurred}:
\ES{, ingresamos en la consola algo que definitivamente no es un número, como \q{asdasd}. \scanf finaliza con 0 en \EAX, lo que indica que ha ocurrido un error:}

\begin{figure}[H]
\centering
\includegraphics[scale=\FigScale]{patterns/04_scanf/3_checking_retval/olly_2.png}
\caption{\olly: \scanf \RU{закончился с ошибкой}\EN{returning error}\ES{devolviendo error}}
\label{fig:scanf_ex3_olly_2}
\end{figure}

\RU{Вместе с этим мы можем посмотреть на локальную переменную в стеке~--- она не изменилась.}
\EN{We can also check the local variable in the stack and note that it has not changed.}
\ES{También podemos verificar la variable local en la pila y notar que no ha cambiado.}
\RU{Действительно, ведь что туда записала бы функция \scanf}\EN{Indeed, what would \scanf write there}?
\ES{De hecho, ¿qué escribiría \scanf allí}?
\RU{Она не делала ничего кроме возвращения нуля}\EN{It simply did nothing except returning zero}.
\ES{Simplemente no hizo nada excepto devolver cero.}

\RU{Попробуем ещё немного \q{хакнуть} нашу программу}\EN{Let's try to \q{hack} our program}.
\RU{Щелкнем правой кнопкой на}\EN{Right-click on}\ES{Intentemos "hackear" un poco nuestro programa. Haga clic derecho en} \EAX, 
\RU{там, в числе опций, будет также}\EN{Among the options there is}\ES{, entre las opciones está} \q{Set to 1}.
\RU{Это нам и нужно}\EN{This is what we need}\ES{. Esto es lo que necesitamos.}

\RU{В \EAX теперь 1, последующая проверка пройдет как надо, и \printf выведет значение переменной
из стека.}
\EN{We now have 1 in \EAX, so the following check is to be executed as intended, 
and \printf will print the value of the variable in the stack.}
\ES{Ahora tenemos 1 en \EAX, por lo que la siguiente comprobación se ejecutará como se esperaba, y \printf imprimirá el valor de la variable en la pila.}

\RU{Запускаем (F9) и видим в консоли следующее:}
\EN{When we run the program (F9) we can see the following in the console window:}
\ES{Al ejecutar el programa (F9) podemos ver lo siguiente en la ventana de la consola:}

\begin{figure}[H]
\centering
\includegraphics[scale=\FigScale]{patterns/04_scanf/3_checking_retval/olly_3.png}
\caption{\RU{консоль}\EN{console window}\ES{ventana de la consola}}
\end{figure}

\RU{Действительно}\EN{Indeed}\ES{De hecho}, 1850296084 \RU{это десятичное представление числа в стеке}
\EN{is a decimal representation of the number in the stack}
\ES{es la representación decimal del número en la pila} (\TT{0x6E494714})!
